\chapter{Infinitud, enumerabilitat i seqüències}\label{cap:infinitud}

\begin{objectius}
\apartat{Prerequisits} Del capítol~\ref{cap:finits}: finitud interna, cardinal, principi dels calaixos, l'equivalència mono/epi/iso en cardinal igual i la tècnica de codificar aplicacions entre objectes variables dins un exponencial fix. Del capítol~\ref{cap:naturals}: recursió (simple i amb paràmetres), inducció, aritmètica i bon ordre de \(\N\). Del capítol~\ref{cap:parts-interna}: quantificadors interns, per tallar parts definides per condicions.
\apartat{Objectius} En acabar el capítol, el lector hauria de poder:
\begin{itemize}
  \item distingir \enquote{infinit} (no finit) d'\enquote{infinit de Dedekind} (endomonomorfisme no exhaustiu), demostrar que \(\N\) és infinit de Dedekind i que Dedekind implica infinit sense cap axioma;
  \item demostrar que infinit implica Dedekind amb l'elecció numerable, construint primer una família de llistes injectives i després, sense més elecció, una successió injectiva;
  \item definir enumerable i numerable, i demostrar l'estabilitat per imatge, part, suma i producte sense elecció, amb un aparellament \(\N\cong\N\times\N\) construït per recursió;
  \item demostrar que tota part de \(\N\) és enumerable (el mínim substitueix l'elecció) i que tot enumerable infinit és isomorf a \(\N\);
  \item construir \(\operatorname{List}(A)\) com a part d'un exponencial fix i demostrar-ne el principi de recursió amb el truc de la diagonal sobre la llargada;
  \item situar l'elecció numerable: on cal (unió numerable d'enumerables, dues infinituds) i on no (parts de \(\N\), llistes).
\end{itemize}
\end{objectius}

El capítol anterior va internalitzar la finitud: un
conjunt és finit quan és isomorf a algun segment inicial \(\seg{n}\)
(definició~\ref{def:conjunt-finit}), i el seu tret distintiu és
l'equivalència, en cardinal igual, entre injectivitat, exhaustivitat i
bijectivitat (teorema~\ref{teo:endo-finit}). El present capítol pren
aquesta darrera equivalència com a frontissa: un conjunt és \emph{infinit}
quan simplement no és finit (definició~\ref{def:conjunt-infinit}), i el que
la fa manejable és que en un conjunt infinit aquella equivalència falla, de
manera que admet un monomorfisme estricte en si mateix. Aquesta darrera
propietat ---la infinitud de \emph{Dedekind}--- no és la definició
d'infinit, sinó una noció a priori distinta; la seva equivalència amb «no
finit» depèn de l'elecció, que explicitarem més endavant
(teorema~\ref{teo:infinit-implica-dedekind}). A partir d'aquí desenvolupem la
comparació de cardinals \emph{sense} la restricció de finitud ---els
conjunts enumerables, els subobjectes de \(\N\), les llistes finites i
les unions indexades numerablement--- amb un fil conductor constant:
distingir en cada pas què es demostra amb els axiomes ja assumits i què en
requereix un de nou. L'axioma nou que reapareixerà és l'\emph{elecció
numerable}; el senyalarem cada cop que intervingui.

\section{Conjunts infinits i infinitud de Dedekind}\label{sec:infinits}

La manera més directa de dir que un conjunt és infinit és negar que sigui
finit. Aquesta definició, però, és purament externa a l'estructura del
conjunt: no exhibeix cap testimoni de la infinitud. Dedekind en va donar
una de positiva ---un conjunt és infinit si es pot posar en bijecció amb
una part pròpia--- que sí que n'aporta un, i que resulta molt més
manejable. Introduïm totes dues i n'estudiem la relació.

\begin{definicio}[Conjunt infinit]\index{Conjunt infinit}\label{def:conjunt-infinit}
Un conjunt \(A\) és \emph{infinit} si no és finit; és a dir, si
\(\operatorname{Card}(A)\) és inicial (definició~\ref{def:conjunt-finit}):
no hi ha cap punt \emph{intern} \(n\colon\mathbf1\to\N\) amb
\(\seg{n}\approx A\).
\end{definicio}

\begin{observacio}[Lectura externa amb numerals estàndard]\label{obs:infinit-numeral}
En particular, si \(A\) és infinit, aleshores no és isomorf a cap ordinal
\([k]\) associat a un numeral estàndard \(k\): si ho fos, la coincidència
\(\seg{\bar k}\cong[k]\) (observació~\ref{obs:segment-numeral}) donaria
\(\bar k\in\operatorname{Card}(A)\), contra la definició. El recíproc, en
canvi, no val en general: un punt global intern \(n\colon\mathbf1\to\N\)
no ha de ser necessàriament un numeral estàndard, de manera que
\(\operatorname{Card}(A)\) pot ser no inicial ---i \(A\) finit en el sentit
intern--- sense que \(A\) sigui isomorf a cap \([k]\) estàndard. Només es
pot afirmar el recíproc sota una hipòtesi externa que garanteixi que el
testimoni intern de finitud és estàndard.
\end{observacio}

\begin{definicio}[Infinit de Dedekind]\index{Infinit de Dedekind}\label{def:dedekind-infinit}
Un conjunt \(A\) és \emph{infinit de Dedekind} si existeix un monomorfisme
\(f\colon A\hookrightarrow A\) que no és exhaustiu; equivalentment, si
existeix un monomorfisme \(A\hookrightarrow A\) la imatge del qual és una
part pròpia de \(A\).
\end{definicio}

\begin{observacio}[Les dues nocions i el cas finit]\label{obs:dues-infinituds}
Sobre un conjunt finit les dues definicions coincideixen amb la negació
l'una de l'altra: pel teorema~\ref{teo:endo-finit}, un monomorfisme d'un
conjunt finit en si mateix és sempre exhaustiu, de manera que cap conjunt
finit és infinit de Dedekind. La qüestió delicada és la recíproca ---si
tot conjunt infinit és infinit de Dedekind---, i és precisament on
apareixerà un axioma addicional. Convé, doncs, mantenir les dues nocions
separades fins que en tinguem la relació exacta.
\end{observacio}

El primer fet, i el que justifica tota la teoria, és que \(\N\) mateix és
infinit de Dedekind, i el testimoni és el successor.

\begin{proposicio}[\(\N\) és infinit de Dedekind]\label{prop:N-dedekind}
El successor \(s\colon\N\to\N\) és un monomorfisme no exhaustiu. Per tant
\(\N\) és infinit de Dedekind, i en conseqüència infinit.
\end{proposicio}

\begin{prova}
Que \(s\) és un monomorfisme és el teorema~\ref{teo:peano}(1): de
\(s\circ m=s\circ n\) se segueix \(m=n\). Que no és exhaustiu és el
teorema~\ref{teo:peano}(2): \(0\) no és de la forma \(s\circ n\)
per a cap \(n\), de manera que el punt \(0\colon\mathbf1\to\N\) no factoritza
per \(s\), i la imatge de \(s\) és la part pròpia complementària del zero.
Així \(s\) exhibeix \(\N\) com a infinit de Dedekind. Que \(\N\) és
infinit se segueix de l'observació~\ref{obs:dues-infinituds}: si fos
finit, \(s\) hauria de ser exhaustiu pel teorema~\ref{teo:endo-finit},
contra el que acabem de veure.
\end{prova}

Una direcció de l'equivalència entre les dues infinituds no costa res i
no fa servir cap axioma nou: qui és infinit de Dedekind és infinit.

\begin{proposicio}[Dedekind implica infinit]\label{prop:dedekind-implica-infinit}
Tot conjunt infinit de Dedekind és infinit. La implicació no usa cap
axioma més enllà dels ja assumits.
\end{proposicio}

\begin{prova}
És el contrarecíproc de la primera part de
l'observació~\ref{obs:dues-infinituds}. Si \(A\) fos finit i infinit de
Dedekind alhora, tindríem un monomorfisme \(A\hookrightarrow A\) no
exhaustiu; però pel teorema~\ref{teo:endo-finit} tot monomorfisme d'un
conjunt finit en si mateix és exhaustiu, una contradicció. Per tant un
conjunt infinit de Dedekind no pot ser finit.
\end{prova}

La direcció oposada ---tot conjunt infinit és infinit de Dedekind--- és
la que requereix construir, dins d'un conjunt del qual només sabem que no
és finit, una còpia de \(\N\); i triar-la exigeix una successió d'eleccions.

\begin{teorema}[Infinit implica Dedekind, amb elecció numerable]\label{teo:infinit-implica-dedekind}
Suposem l'\emph{axioma d'elecció numerable}
(axioma~\ref{ax:eleccio-numerable}). Aleshores tot conjunt infinit és
infinit de Dedekind. Equivalentment, sota aquest axioma un conjunt és
infinit si i només si admet un monomorfisme \(\N\hookrightarrow A\).
\end{teorema}

\begin{prova}
Sigui \(A\) infinit. La dificultat és que no podem anar \emph{triant}
\(a_{n+1}\) dins el complement \(A\setminus\{a_0,\dots,a_n\}\): aquest
complement depèn de les tries anteriors, de manera que la família de
conjunts no buits sobre la qual triaríem no està donada per endavant i
l'axioma~\ref{ax:eleccio-numerable} no s'hi aplica directament. Ho evitem
fixant \emph{abans} de triar una única família numerable, independent de
tota elecció.

\emph{La família.} Per a cada \(n\colon\mathbf1\to\N\), sigui
\[
I_n\ :=\ \operatorname{Inj}\bigl(\seg{s\circ n},A\bigr)\ \hookrightarrow\
A^{\seg{s\circ n}}
\]
la part de l'exponencial \emph{fix} \(A^{\seg{s\circ n}}\) formada pels
monomorfismes \(\seg{s\circ n}\hookrightarrow A\) ---les llistes
injectives de llargada \(n+1\)---, tallada per la condició interna
d'injectivitat (quantificadors acotats del capítol~\ref{cap:parts-interna}).
Cada \(I_n\) és \emph{no inicial}, és a dir, per a cada \(n\) existeix un
monomorfisme \(\seg{s\circ n}\hookrightarrow A\). Ho veiem per inducció
sobre \(n\), sense elecció. Per a \(n=0\), \(A\) és no inicial ---no és
finit, en particular no és buit---, de manera que hi ha un punt
\(\mathbf1\cong\seg{1}\hookrightarrow A\). Suposem un monomorfisme
\(g\colon\seg{s\circ n}\hookrightarrow A\); la seva imatge és una part
finita de \(A\) (cardinal \(n+1\)), i com que \(A\) no és finit, aquesta
imatge no és tot \(A\) (altrament \(\seg{s\circ n}\approx A\) i \(A\) seria
finit, definició~\ref{def:conjunt-finit}). El seu complement és, doncs, no
inicial i en conté algun punt \(x\colon\mathbf1\to A\); estenent \(g\) amb
\(x\) sobre el sumand nou de \(\seg{s\circ(s\circ n)}\cong\seg{s\circ
n}+\mathbf1\) (proposició~\ref{prop:segment-successor}) s'obté un
monomorfisme \(\seg{s\circ(s\circ n)}\hookrightarrow A\), perquè \(x\) no
és a la imatge de \(g\). Això tanca la inducció i prova que
\(I_n\) és no inicial per a tot \(n\). Noteu que aquesta inducció només
afirma \emph{existència} fibra a fibra: no proporciona una tria uniforme de
les \(g\), i és per obtenir-ne una que caldrà l'elecció numerable.

Aquestes parts es reuneixen en una \emph{sola} família interna. Codifiquem
cada \(f\in I_n\), com a la definició~\ref{def:card-A}, per la seva
extensió \(\bar f\colon\N\to A+\mathbf1\) que val \(f\) sobre
\(\seg{s\circ n}\) i el punt afegit fora; així totes les llistes
injectives, de qualsevol llargada, viuen dins l'exponencial \emph{fix}
\((A+\mathbf1)^{\N}\). Sigui
\[
E\ :=\ \bigl\{\,(n,\bar f)\ \mid\ n\colon\mathbf1\to\N,\ f\in I_n\,\bigr\}\
\hookrightarrow\ \N\times(A+\mathbf1)^{\N},
\]
la part corresponent (tallada per equalitzadors i condicions acotades), i
\(p\colon E\to\N\) la projecció, la fibra de la qual sobre \(n\) és
exactament \(I_n\). Com que cada \(I_n\) és no inicial, tot punt \(n\) té preimatge per \(p\),
i \(p\) és un epimorfisme (exhaustivitat sobre punts,
proposició~\ref{prop:epi-exhaustiu-global}), regular pel
corol·lari~\ref{cor:epi-regular-sense-eleccio}.
Aquesta família \emph{ja existeix} abans de qualsevol tria.

\emph{La tria.} L'axioma d'elecció numerable
(axioma~\ref{ax:eleccio-numerable}) dona una secció
\(\sigma\colon\N\to E\) de \(p\); és a dir, per a cada \(n\) un monomorfisme
\(f_n\colon\seg{s\circ n}\hookrightarrow A\) escollit \emph{d'un sol cop}
per a tota la família alhora, i no pas un a un. Escrivim
\(f_n(k)\in A\) per al valor de la \(k\)-èsima component,
\(k<n+1\).

\emph{El monomorfisme \(\N\hookrightarrow A\).} A partir de la successió
\((f_n)_n\) construïm \(a\colon\N\to A\) \emph{sense més elecció}. La
temptació és definir \(a_n\) com el primer valor de \(f_n\) que no és cap
dels \(a_0,\dots,a_{n-1}\); però això és una recursió que depèn de tota la
història anterior, i no té la forma de l'axioma~\ref{ax:naturals}. Ho
evitem fent servir llistes prou llargues perquè els valors ja usats no
puguin esgotar-les. Per a cada punt \(n\), sigui
\[
V_n\ \hookrightarrow\ A
\]
la part dels \(x\) tals que \(x=f_{2^m-1}(i)\) per a algun \(m<n\) i algun
\(i<2^m\): és la imatge de la part de \(\N\times\N\) dels parells
\((m,i)\) amb \(m<n\), \(i<2^m\) per l'aplicació
\((m,i)\mapsto f_{2^m-1}(i)\), i la regla \(n\mapsto V_n\) és una
aplicació \(\N\to P(A)\) (quantificadors interns,
capítol~\ref{cap:parts-interna}); aquí \(2^m\) és l'exponenciació de
naturals (definició~\ref{def:producte-exp}) i \(2^m-1\) el seu
predecessor. La part \(V_n\) és finita i \(|V_n|<2^n\). En efecte, el conjunt
d'índexs \(J_n:=\{(m,i):m<n,\ i<2^m\}\) és finit amb \(|J_n|+1=2^n\): per
inducció sobre \(n\), \(|J_0|+1=0+1=2^0\), i
\(J_{n+1}\cong J_n+\seg{2^n}\) (s'hi afegeix la fila \(m=n\)), de manera que,
pel recompte de la suma (lema~\ref{lema:recompte-suma}),
\(|J_{n+1}|+1=(|J_n|+1)+2^n=2^n+2^n=2^{n+1}\). A més, \(V_n\) és imatge de
\(J_n\): un epimorfisme \(J_n\twoheadrightarrow V_n\) té una secció ---enviant
cada punt al mínim índex que hi va a parar, ja que \(J_n\) és ben ordenat
com a part finita de \(\N\times\N\) ordenada lexicogràficament, o bé per
l'elecció---, que és un monomorfisme \(V_n\hookrightarrow J_n\); pel
lema~\ref{lema:part-finita}, \(V_n\) és finita amb
\(|V_n|\le|J_n|<2^n\). En canvi, la llista
\(f_{2^n-1}\) té \(2^n\) valors distints. Pel principi dels calaixos
(teorema~\ref{teo:calaixos}) algun valor de \(f_{2^n-1}\) no és a \(V_n\).
Definim
\[
a_n:=f_{2^n-1}(k_n),
\qquad
k_n:=\min\{\,k<2^n\ :\ f_{2^n-1}(k)\notin V_n\,\},
\]
on el mínim existeix pel bon ordre de \(\N\)
(teorema~\ref{teo:bon-ordre-N}) i és únic: \(a\) és una aplicació
\(\N\to A\) definida sense cap tria (internament, el seu graf és la part
de \(\N\times A\) tallada per la condició, funcional i total). Per
construcció \(a_n\notin V_n\), mentre que \(a_m\in V_n\) per a tot
\(m<n\), perquè \(a_m\) és un valor de \(f_{2^m-1}\); per tant
\(a_n\ne a_m\) per a \(m<n\), i \(a\) és injectiva sobre punts, és a dir,
un monomorfisme (separació global).

Finalment, sigui \(C\) el complement de la imatge de \(a\). El primer
isomorfisme de
\[
A\;\cong\;\N + C\;\xrightarrow{\;s+\id_C\;}\;\N+C\;\cong\;A
\]
descompon \(A\) en aquesta imatge (isomorfa a \(\N\) via \(a\)) i \(C\); la
composició reprodueix el desplaçament del successor en el sumand \(\N\) i
omet el punt corresponent a \(0\), de manera que és un monomorfisme
\(A\hookrightarrow A\) no exhaustiu. Així \(A\) és infinit de Dedekind. La
reformulació és immediata: \(a\colon\N\hookrightarrow A\) és el testimoni, i
recíprocament un tal monomorfisme dona, transportant el successor, un
endomorfisme no exhaustiu de \(A\).
\end{prova}

\begin{observacio}[Per què cal l'elecció aquí]\label{obs:necessitat-eleccio-dedekind}
La prova anterior no pot evitar l'elecció numerable amb els axiomes que
tenim. Fixeu-vos on entra i on no. La família \((I_n)_n\) dels monomorfismes
\(\seg{s\circ n}\hookrightarrow A\) està definida \emph{abans} de tota
elecció i cada \(I_n\) és no inicial; això és una afirmació d'existència
fibra a fibra, provada per inducció sense triar res. El que \emph{no}
tenim, sense elecció, és una \emph{secció} uniforme d'aquesta família ---una
tria simultània d'un \(f_n\) a cada fibra---: en cada \(I_n\) hi ha molts
monomorfismes i res no en distingeix cap de canònic. És exactament aquest
pas, i només aquest, el que consumeix l'axioma~\ref{ax:eleccio-numerable};
un cop tenim la successió \((f_n)_n\), la resta de la construcció (el pas al
mínim índex nou dins una llista prou llarga, el desplaçament del
successor) torna a ser canònica. Es
tracta d'un ús genuí i, en el marc habitual, inevitable: sense cap forma
d'elecció és consistent que existeixin conjunts infinits ---no isomorfs a
cap \([k]\) estàndard--- que \emph{no} siguin infinits de Dedekind, els
anomenats conjunts infinits \emph{amb finitud de Dedekind}. La direcció
fàcil (proposició~\ref{prop:dedekind-implica-infinit}), en canvi, no
necessita res. Aquesta asimetria és el motiu pel qual mantenim les dues
nocions d'infinitud amb noms diferents.
\end{observacio}

Enunciem l'axioma que hem fet servir, en la forma que ens caldrà després.

\begin{axioma}[Elecció numerable]\label{ax:eleccio-numerable}
Sigui \(p\colon E\twoheadrightarrow\N\) un epimorfisme regular ---una
família de conjunts no buits indexada per \(\N\), on la fibra sobre cada
\(n\) és la part no buida corresponent. Aleshores \(p\) admet una secció
\(\sigma\colon\N\to E\), és a dir, una aplicació amb \(p\circ\sigma=\id_\N\)
que tria un element de cada fibra.
\end{axioma}

\begin{observacio}[Elecció, elecció numerable i una nota terminològica]\label{obs:eleccio-global-vs-numerable}
L'\emph{axioma d'elecció} del llibre (capítol~\ref{cap:relacions}) diu que
\emph{tot} epimorfisme regular s'escindeix; l'elecció numerable
(axioma~\ref{ax:eleccio-numerable}) n'és el cas particular en què l'índex
és \(\N\), i per tant se'n dedueix. Si el llibre assumeix l'elecció plena,
l'axioma~\ref{ax:eleccio-numerable} no és una hipòtesi nova sinó una
instància de la que ja hi ha, i el
teorema~\ref{teo:infinit-implica-dedekind} és un teorema sense reserves.
L'enunciem a part perquè la major part de la teoria d'aquest capítol només
en necessita la versió feble, i perquè convé veure exactament on entra
alguna forma d'elecció: en tot el que segueix, els resultats que no
l'esmenten no en depenen.

Una advertència de vocabulari, que serà important en els capítols
transfinits. Anomenem aquest principi simplement \emph{axioma d'elecció};
evitem el terme «elecció global», que en teoria de conjunts i de classes
designa un principi més fort ---una funció d'elecció única sobre tota la
classe dels conjunts--- que aquí ni enunciem ni necessitem. Igualment, cal
no confondre l'elecció amb cap principi de \emph{substitució} o de reunió de
famílies: l'elecció tria elements dins les fibres d'una família \emph{que ja
existeix} (com \(p\colon E\to\N\) a la prova anterior), mentre que la
substitució serviria per garantir l'existència d'una família d'objectes
definida uniformement. Són funcions lògiques distintes, i en aquest capítol
només usem la primera.
\end{observacio}

\section{Conjunts enumerables i numerables}\label{sec:enumerables}

Un conjunt és enumerable quan els seus elements es poden \emph{llistar}
amb els naturals, potser amb repeticions. La manera estructural de dir-ho
és que hi ha una aplicació exhaustiva des de \(\N\).

\begin{definicio}[Conjunt enumerable]\index{Conjunt enumerable}\label{def:enumerable}
Un conjunt \(A\) és \emph{enumerable} si és inicial o bé existeix un
epimorfisme regular \(e\colon\N\twoheadrightarrow A\). Una tal \(e\) és una
\emph{enumeració} de \(A\): l'element \(a_n:=e\circ n\) és el terme
\(n\)-èsim, i l'exhaustivitat diu que tot element generalitzat de \(A\)
és, localment, algun \(a_n\).
\end{definicio}

\begin{definicio}[Conjunt numerable]\index{Conjunt numerable}\label{def:numerable}
Un conjunt és \emph{numerable} si és finit
(definició~\ref{def:conjunt-finit}) o enumerable infinit. Un conjunt
enumerable és, doncs, o bé finit o bé numerable infinit.
\end{definicio}

\begin{observacio}[Enumerable, sense elecció]\label{obs:enumerable-sense-eleccio}
La definició~\ref{def:enumerable} \emph{no} amaga cap elecció: demana una
única aplicació \(e\), no una successió de tries. Per això tota la
combinatòria d'enumerabilitat que en depèn ---imatges d'enumerables,
unions de dues enumeracions, producte de dues--- es fa sense cap axioma
nou. L'elecció numerable només reapareixerà quan calgui triar \emph{una
enumeració per a cadascun} d'una família infinita de conjunts, no quan en
tinguem una de sola.
\end{observacio}

Les operacions bàsiques preserven l'enumerabilitat. Reunim els fets
elementals, tots demostrables sense elecció.

\begin{proposicio}[Estabilitat de l'enumerabilitat]\label{prop:estabilitat-enumerable}
Siguin \(A\) i \(B\) enumerables.
\begin{enumerate}
\item Tota imatge d'\(A\) és enumerable: si \(g\colon A\to C\) és
exhaustiva, \(C\) és enumerable.
\item Tota part d'\(A\) és enumerable.
\item La suma \(A+B\) és enumerable.
\item El producte \(A\times B\) és enumerable.
\end{enumerate}
Cap dels quatre enunciats no requereix elecció.
\end{proposicio}

\begin{prova}
(1) Si \(e\colon\N\twoheadrightarrow A\) és una enumeració i \(g\) és
exhaustiva, la composició \(g\circ e\colon\N\twoheadrightarrow C\) ho és
(l'exhaustivitat es preserva per composició, capítol~\ref{cap:imatges}),
i és una enumeració de \(C\); el cas inicial és trivial.

(2) Sigui \(m\colon P\hookrightarrow A\) una part no inicial. Fixat un punt
\(p_0\colon\mathbf1\to P\) ---que existeix perquè \(P\) no és inicial---,
definim \(r\colon A\to P\) enviant els elements de \(P\) en si mateixos i
la resta a \(p_0\); això és una retracció de \(m\), construïda amb el cas
finit de decisió del complement (capítol~\ref{cap:valors}) i sense cap
tria infinita. Aleshores \(r\circ e\colon\N\twoheadrightarrow P\) és una
enumeració. El cas \(P\) inicial és immediat.

(3) De dues enumeracions \(e\colon\N\twoheadrightarrow A\) i
\(f\colon\N\twoheadrightarrow B\) i d'un isomorfisme
\(\N\cong\N+\N\) ---per exemple la partició en parells i senars, obtinguda
de l'aritmètica del capítol~\ref{cap:naturals}--- s'obté
\(\N\cong\N+\N\xrightarrow{e+f}A+B\), exhaustiva.

(4) N'hi ha prou amb una bijecció \(\N\cong\N\times\N\) ---un
\emph{aparellament}--- perquè aleshores
\(\N\cong\N\times\N\xrightarrow{e\times f}A\times B\) és exhaustiva. Un tal
aparellament es construeix explícitament dins la teoria \emph{sense divisió},
enumerant \(\N\times\N\) per diagonals: es recorre primer la diagonal
\(i+j=0\), després \(i+j=1\), i així successivament, i dins de cada diagonal
s'ordena per \(j\) creixent. Formalment, sigui \(\nu\colon\N\times\N\to\N\times\N\) l'aplicació
\enquote{següent parell en l'ordre per diagonals}, definida per casos
sobre el punt \(i\) (zero o successor, lema~\ref{lema:zero-o-successor}):
\(\nu(0,j)=(s\circ j,0)\) i \(\nu(s\circ i',j)=(i',s\circ j)\). La recursora
\(u\colon\N\to\N\times\N\) del sistema \((\N\times\N,\langle0,0\rangle,\nu)\)
recorre les diagonals, i és bijectiva: injectiva perquè \(\nu\) ho és i
\((0,0)\) no és a la seva imatge (el mateix argument que a les equacions
de Peano), i exhaustiva perquè la part dels \((i,j)\) assolits conté
\((0,0)\) i és tancada per \(\nu\), i tot parell s'obté de \((0,0)\)
iterant \(\nu\) (inducció sobre \(i+j\) i després sobre \(j\)). Els detalls
són un exercici de la part pràctica. No cal cap divisió ni la fórmula
tancada \(\tfrac12(i+j)(i+j+1)+j\), que és correcta però pressuposa
aritmètica que el llibre no ha desenvolupat. És, en tot cas, una
aplicació concreta, no una tria.
\end{prova}

\begin{observacio}[Lectura clàssica]\label{obs:enumerable-classica}
En termes clàssics, un conjunt no buit és enumerable exactament quan els
seus elements es poden disposar en una successió
\(a_0,a_1,a_2,\dots\) que els reculli tots ---amb repeticions permeses.
La proposició~\ref{prop:estabilitat-enumerable} recull els fets habituals:
la imatge d'un conjunt enumerable, qualsevol dels seus subconjunts, i la
unió i el producte de dos enumerables són enumerables; en particular
\(\Z\) i \(\Q\) ho són, com es comprova aparellant-los amb \(\N\). El punt
que la presentació estructural fa visible, i que sovint queda amagat, és
que aquests resultats no fan servir cap elecció: totes les enumeracions
involucrades es construeixen a partir d'una o dues enumeracions donades,
mitjançant aparellaments explícits.
\end{observacio}

\section{\texorpdfstring{Subobjectes de \(\N\)}{Subobjectes de N}}\label{sec:subobjectes-N}

Els subobjectes de \(\N\) ---les seves parts--- són l'exemple bàsic i
mereixen un tractament propi, perquè hi funciona una construcció que en
general falla: es pot triar sempre el \emph{mínim}. Això fa que tot
subobjecte de \(\N\) sigui enumerable sense recórrer a cap elecció.

\begin{teorema}[Els subobjectes de \(\N\) són enumerables]\label{teo:subobjectes-N}
Tota part \(P\hookrightarrow\N\) és enumerable, i la seva enumeració es
construeix sense elecció. A més, \(P\) és finit o bé isomorf a \(\N\); en
el segon cas, l'enumeració es pot prendre estrictament creixent i
bijectiva.
\end{teorema}

\begin{prova}
Si \(P\) és inicial, és enumerable per definició. Suposem \(P\) no
inicial. La clau és el \emph{principi del mínim}: tota part no inicial de
\(\N\) té un element mínim, la qual cosa és una reformulació de la
inducció (teorema~\ref{teo:induccio}) i no usa elecció, perquè el mínim,
quan existeix, és \emph{únic} i per tant canònic. Definim una autoaplicació \(t\colon P\to P\): \(t(x)\) és el mínim dels
elements de \(P\) estrictament més grans que \(x\) si n'hi ha, i \(x\)
altrament. És una aplicació genuïna, perquè el seu graf ---la part de
\(P\times P\) dels \((x,y)\) amb \(y\in P\), \(x<y\) i
\(\forall z\in P\,(x<z\Rightarrow y\le z)\), unida a la part dels \((x,x)\)
amb \(\neg\exists z\in P\,(x<z)\)--- és funcional (el mínim és únic) i
total (o hi ha elements més grans o no), i es talla amb els
quantificadors interns. La successió \(a\colon\N\to P\) és la recursora
del sistema \((P,\min P,t)\): \(a_0\) és el mínim de \(P\) i
\(a_{n+1}=t(a_n)\). És una recursió simple de l'axioma~\ref{ax:naturals},
sense cap tria.

Dos casos. Si per a algun \(n\) no queda cap element de \(P\) més gran que
\(a_n\), la successió s'estabilitza i la seva imatge
\(\{a_0<a_1<\dots<a_n\}\) és exactament \(P\), que és finit. Si en cada pas
sempre en queda un de més gran, \(a\) és estrictament creixent, per tant
un monomorfisme, i la seva imatge és tot \(P\) perquè cap element de \(P\)
pot ser omès. En efecte, \(a_i\ge i\) per a tot \(i\) (inducció: \(a_{i+1}>a_i\)),
de manera que, donat \(k\in P\), la part dels \(i\) amb \(k<a_i\) és no
inicial (conté \(s\circ k\)) i té un mínim \(j\); com que \(a_0=\min P\le k\),
\(j=i+1\) per a algun \(i\), amb \(a_i\le k<a_{i+1}\). Si fos \(a_i<k\), aleshores
\(k\) seria un element de \(P\) estrictament més gran que \(a_i\) i menor que
\(a_{i+1}=t(a_i)\), contra la definició de \(t\) com a mínim d'aquests
elements; per tant \(k=a_i\). Aleshores
\(a\colon\N\to P\) és una bijecció creixent, i \(P\cong\N\).
\end{prova}

\begin{corollari}[Caracterització dels enumerables infinits]\label{cor:enumerable-infinit-N}
Un conjunt és enumerable infinit si i només si és isomorf a \(\N\). En
conseqüència, tots els conjunts numerables infinits són isomorfs entre si,
i \(\N\) n'és el representant.
\end{corollari}

\begin{prova}
Si \(A\cong\N\), és clarament enumerable (per la identitat) i infinit
(proposició~\ref{prop:N-dedekind}). Recíprocament, sigui \(A\) enumerable
infinit i \(e\colon\N\twoheadrightarrow A\) una enumeració. Volem un
monomorfisme \(A\hookrightarrow\N\); un cop el tinguem, la seva imatge és
una part de \(\N\) isomorfa a \(A\), que pel teorema~\ref{teo:subobjectes-N}
és finita o isomorfa a \(\N\), i com que \(A\) és infinit ha de ser el
segon cas. Per construir el monomorfisme sense elecció, definim
\(r\colon A\to\N\) enviant cada element \(a\) al \emph{mínim} índex \(n\)
amb \(e_n=a\); aquest mínim existeix i és únic pel principi del mínim, de
manera que \(r\) està ben definida i és una secció de \(e\), en particular
un monomorfisme. Per tant \(A\cong\N\).
\end{prova}

\begin{observacio}[El mínim substitueix l'elecció]\label{obs:minim-vs-eleccio}
Els dos resultats anteriors s'obtenen sense elecció \emph{precisament}
perquè a \(\N\) hi ha una tria canònica ---el mínim--- disponible a cada
pas. És el mateix fenomen que fa que l'elecció sigui innecessària sobre
conjunts ben ordenats: quan cada part no buida té un element distingit, no
cal cap axioma per triar-lo. En un conjunt infinit qualsevol, sense un tal
ordre, aquesta comoditat desapareix, i és per això que el
teorema~\ref{teo:infinit-implica-dedekind} sí que necessita l'elecció
numerable mentre que el teorema~\ref{teo:subobjectes-N} no.
\end{observacio}

\section{Llistes finites i iteració}\label{sec:llistes}

Una \emph{llista finita} d'elements d'un conjunt \(A\) és una successió
finita \((a_0,\dots,a_{n-1})\) de llargada arbitrària però determinada.
Reunir totes les llistes de totes les llargades en un sol conjunt és la
construcció que fa d'\(A\) el generador d'un monoide lliure, i és el
prototip d'una unió indexada per \(\N\) que, en aquest cas, es fa sense
elecció.

\begin{definicio}[Conjunt de llistes finites]\index{Conjunt de llistes finites}\label{def:llistes}
El \emph{conjunt de les llistes finites} d'\(A\), denotat
\(\operatorname{List}(A)\), es defineix com una part de l'exponencial
\emph{fix} \((A+\mathbf1)^{\N}\) de les successions amb un valor de
terminació. Escrivim \(i_A\colon A\to A+\mathbf1\) i
\(\ast\colon\mathbf1\to A+\mathbf1\) per a les injeccions. Una successió
\(u\colon\N\to A+\mathbf1\) és una \emph{llista} quan el conjunt de
posicions \emph{ocupades} ---les \(k\) amb \(u(k)=i_A(a)\) per a algun
\(a\)--- és un segment inicial finit; és a dir, quan existeix
\(n\colon\mathbf1\to\N\) tal que
\begin{enumerate}
\item[(l1)] per a tot \(k<n\), \(u(k)=i_A(a)\) per a un únic \(a\in A\);
\item[(l2)] per a tot \(k\) amb \(k\not<n\), \(u(k)=\ast\).
\end{enumerate}
Aquesta \(n\) és única quan existeix (és la \emph{llargada}), i les dues
condicions es tallen amb quantificadors acotats i igualtats sobre
l'exponencial fix (capítol~\ref{cap:parts-interna}), de manera que
\[
\operatorname{List}(A)\ :=\ \bigl\{\,u\in(A+\mathbf1)^{\N}\ \mid\ u\text{
satisfà (l1)--(l2) per a algun }n\,\bigr\}\ \hookrightarrow\
(A+\mathbf1)^{\N}
\]
és una part interna ben definida. La successió constant \(\ast\)
(cas \(n=0\)) és la llista buida
\(\varepsilon\colon\mathbf1\to\operatorname{List}(A)\).
\end{definicio}

\begin{observacio}[La condició de segment inicial és essencial]\label{obs:llistes-segment}
No n'hi ha prou d'exigir que \(u\) sigui \emph{eventualment constant} igual
a \(\ast\): aquesta condició, més feble, admet successions com
\(a,\ast,b,\ast,\ast,\dots\), amb una posició de terminació seguida de nous
elements d'\(A\), que no representen cap llista. La condició (l1)--(l2)
demana que \emph{totes} les posicions ocupades formin un segment inicial
\(\seg{n}\): un cop apareix el primer \(\ast\), tots els termes següents
són \(\ast\). Amb aquesta condició, i només amb ella, la llargada està ben
definida i la lectura
\[
\operatorname{List}(A)\ \cong\ \sum_{n\colon\mathbf1\to\N}A^{\seg{n}}
\;=\;\coprod_{n\in\N}A^{\seg{n}}
\]
és correcta: cada llista de llargada \(n\) es correspon amb la seva
restricció \(\seg{n}\to A\), i \(A^{\seg{0}}\cong\mathbf1\) és la llista
buida. Aquesta igualtat és, doncs, un \emph{teorema} sobre la part just
construïda, no la definició: la suma indexada per \(\N\) del membre dret no
es pot formar directament, perquè els axiomes només donen sumes finites
(capítol~\ref{cap:valors}). Vegeu l'observació~\ref{obs:existencia-llistes}.
\end{observacio}

\begin{observacio}[Sobre l'existència i les dues lectures]\label{obs:existencia-llistes}
La definició~\ref{def:llistes} construeix \(\operatorname{List}(A)\)
\emph{directament}, com una part d'un exponencial fix, i per tant la seva
existència no demana cap suma numerable ni cap principi nou: només
l'exponencial \((A+\mathbf1)^{\N}\) i els quantificadors interns
(capítols~\ref{cap:exponencials} i~\ref{cap:parts-interna}). Aquesta és la
via segura. Les altres dues descripcions habituals ---la suma
\(\coprod_{n}A^{\seg{n}}\) i l'objecte inicial de les
\((\mathbf1+A\times{-})\)-àlgebres--- \emph{no} es prenen aquí com a
definicions: la primera perquè la suma numerable no és una construcció
finita (observació~\ref{obs:llistes-segment}), i la segona perquè afirmar
que aquesta àlgebra inicial \emph{existeix} és precisament el que s'ha de
demostrar. Totes dues es recuperen com a \emph{propietats} de la part ja
construïda: la lectura com a suma és l'observació~\ref{obs:llistes-segment},
i la inicialitat és el contingut de la
proposició~\ref{prop:recursio-llistes}, que provem tot seguit sobre
\(\operatorname{List}(A)\) equipat amb \(\varepsilon\) i \(\operatorname{cons}\).
\end{observacio}

\begin{proposicio}[Recursió i iteració sobre llistes]\label{prop:recursio-llistes}
El conjunt \(\operatorname{List}(A)\) ve equipat amb el punt
\(\varepsilon\colon\mathbf1\to\operatorname{List}(A)\) de la llista buida i
amb l'operació de prefixar un element,
\(\operatorname{cons}\colon A\times\operatorname{List}(A)\to
\operatorname{List}(A)\), i satisfà el principi de recursió corresponent:
per a tot conjunt \(B\) amb un punt \(b\colon\mathbf1\to B\) i una aplicació
\(t\colon A\times B\to B\), existeix una única
\(h\colon\operatorname{List}(A)\to B\) amb
\[
h\circ\varepsilon=b
\qquad\text{i}\qquad
h\circ\operatorname{cons}=t\circ(\id_A\times h).
\]
En particular, la llargada
\(\ell\colon\operatorname{List}(A)\to\N\), la concatenació i la iteració
d'una operació al llarg d'una llista es defineixen totes per aquesta
recursió, sense cap axioma nou.
\end{proposicio}

\begin{prova}
\emph{L'operació \(\operatorname{cons}\).} Sobre la representació de la
definició~\ref{def:llistes}, prefixar \(a\in A\) a una llista \(u\) de
llargada \(n\) és la successió \(\operatorname{cons}(a,u)\) de llargada
\(n+1\) que val \(i_A(a)\) a la posició \(0\) i \(u(k-1)\) ---és a dir, el
desplaçament de \(u\)--- a la posició \(k\) per a \(1\le k<n+1\), i \(\ast\)
a la resta. El desplaçament que reindexa el terme de la posició \(k\) cap a
la posició \(k+1\), i la reassignació de la posició \(0\), són aplicacions
internes sobre
\((A+\mathbf1)^{\N}\) (avaluació, predecessor acotat i definició per casos,
capítols~\ref{cap:exponencials} i~\ref{cap:parts-interna}), i preserven la
condició (l1)--(l2); per tant
\(\operatorname{cons}\colon A\times\operatorname{List}(A)\to
\operatorname{List}(A)\) està ben definida. La llargada compleix
\(\ell(\varepsilon)=0\) i \(\ell(\operatorname{cons}(a,u))=s(\ell(u))\).

\emph{Existència de \(h\).} Escrivim \(L:=\operatorname{List}(A)\). Les
aplicacions \(h_m\colon L\to B\) que \enquote{encerten} sobre les llistes
de llargada \(\le m\) no les construïm amb dominis variables, sinó com a
punts de l'exponencial fix \(B^{L}\), per recursió simple. Sigui
\(\Phi\colon B^{L}\to B^{L}\) la transposada de l'aplicació
\(B^{L}\times L\to B\) definida per casos sobre la descomposició
\(L\cong\mathbf1+A\times L\) (llista buida, o primer element i cua,
recuperats internament, amb \(\operatorname{cons}\) com a isomorfisme):
\[
(h,\varepsilon)\mapsto b,
\qquad
(h,\operatorname{cons}(a,u))\mapsto t\bigl(a,\operatorname{ev}(h,u)\bigr).
\]
Sigui \(H\colon\N\to B^{L}\) la recursora del sistema
\((B^{L},\ulcorner b\circ{!}_L\urcorner,\Phi)\), de manera que
\(H_{0}\) és la constant \(b\) i \(H_{m+1}=\Phi(H_m)\). Per inducció
sobre \(m\) (amb la llargada com a paràmetre), \(H_m\) coincideix amb
\(H_{m'}\) sobre les llistes de llargada \(\le m\) per a tot \(m'\ge m\):
per a \(m=0\), totes valen \(b\) sobre \(\varepsilon\); i si val per a
\(m\), aleshores sobre \(\operatorname{cons}(a,u)\) amb \(\ell(u)\le m\)
es té \(H_{m'+1}(\operatorname{cons}(a,u))=t(a,H_{m'}(u))=t(a,H_m(u))
=H_{m+1}(\operatorname{cons}(a,u))\). Definim \(h\) per la
\emph{diagonal}: \(h:=\operatorname{ev}\circ\langle H\circ\ell,\id_L\rangle\),
és a dir, \(h(u):=H_{\ell(u)}(u)\). Aleshores \(h(\varepsilon)=H_0(\varepsilon)=b\),
i \(h(\operatorname{cons}(a,u))=H_{\ell(u)+1}(\operatorname{cons}(a,u))
=t(a,H_{\ell(u)}(u))=t(a,h(u))\), que són les dues equacions.

\emph{Unicitat.} Si \(h'\) compleix les mateixes dues equacions, aleshores
\(h=h'\) sobre les llistes de llargada \(0\) (totes dues valen \(b\)), i si
coincideixen sobre les de llargada \(\le m\), l'equació del
\(\operatorname{cons}\) força la coincidència sobre les de llargada \(m+1\);
per inducció sobre \(m\) (teorema~\ref{teo:induccio}; la part dels \(m\)
sobre els quals coincideixen és interna, amb un quantificador acotat per
\(\ell\)), \(h=h'\) a tot arreu.
Això estableix la propietat universal ---la inicialitat de
\((\operatorname{List}(A),\varepsilon,\operatorname{cons})\) entre les
\((\mathbf1+A\times{-})\)-àlgebres--- \emph{a partir} de la construcció, i
no a l'inrevés. La llargada s'obté amb \(B=\N\), \(b=0\), \(t=s\circ\pi_2\);
la concatenació amb una llista fixada iterant \(\operatorname{cons}\); i
qualsevol plegat (\emph{fold}) d'una operació binària, amb el \(t\)
corresponent.
\end{prova}

\begin{observacio}[Enumerabilitat de \(\operatorname{List}(A)\)]\label{obs:llistes-enumerables}
Si \(A\) és enumerable, també ho és \(\operatorname{List}(A)\), i sense
elecció. L'argument és en dos passos. Primer, \(\operatorname{List}(\N)\)
és enumerable: la codificació
\(c\colon\operatorname{List}(\N)\to\N\) definida per recursió sobre llistes
(proposició~\ref{prop:recursio-llistes}) amb \(c(\varepsilon)=0\) i
\(c(\operatorname{cons}(k,u))=s\bigl(\pi(k,c(u))\bigr)\), on
\(\pi\colon\N\times\N\to\N\) és un aparellament bijectiu
(proposició~\ref{prop:estabilitat-enumerable}), és bijectiva: injectiva
perquè \(\pi\) i \(s\) ho són i \(0\) no és successor, i exhaustiva per
inducció sobre el codi. Segon, una enumeració \(e\colon\N\twoheadrightarrow A\)
indueix \(\operatorname{List}(e)\colon\operatorname{List}(\N)\to\operatorname{List}(A)\)
(aplicar \(e\) terme a terme, també per recursió sobre llistes), que és
exhaustiva sobre punts: tota llista \((a_0,\dots,a_{n-1})\) és la imatge
de la llista dels mínims índexs \(k_i\) amb \(e(k_i)=a_i\), que es
defineixen sense cap tria. Així
\(\operatorname{List}(e)\circ c^{-1}\colon\N\twoheadrightarrow\operatorname{List}(A)\)
és una enumeració. El que fa possible evitar l'elecció és que les
enumeracions de totes les llistes deriven \emph{uniformement} d'una sola
enumeració d'\(A\) i d'un sol aparellament; és l'observació clau que
separa aquest cas del de la secció següent.
\end{observacio}

\section{Unions indexades numerablement}\label{sec:unions-numerables}

L'últim tema és el que fa explícit el paper de l'elecció numerable en la
seva forma més habitual. Es tracta de reunir una família de conjunts
\((A_n)_{n\in\N}\) en una sola unió i de decidir quan aquesta unió és
enumerable.

\begin{definicio}[Família i unió numerables]\index{Família i unió numerables}\label{def:familia-numerable}
Una \emph{família numerable} de conjunts és una aplicació
\(p\colon E\to\N\); el conjunt \(A_n\) és la fibra
\(p^{-1}(n)\hookrightarrow E\) sobre el punt \(n\)
(capítol~\ref{cap:parts}). La \emph{unió} de la família és la imatge de
\(E\) per una aplicació cap a un conjunt ambient, o simplement \(E\) mateix
quan es pensa com la reunió disjunta \(\coprod_{n\in\N}A_n\); la
\emph{unió indexada} \(\bigcup_{n\in\N}A_n\) dins un conjunt \(U\) que
conté totes les fibres és la imatge (capítol~\ref{cap:imatges}) de \(E\to
U\).
\end{definicio}

\begin{teorema}[Unió numerable d'enumerables, amb elecció numerable]\label{teo:unio-numerable}
Suposem l'axioma d'elecció numerable
(axioma~\ref{ax:eleccio-numerable}). Sigui \((A_n)_{n\in\N}\) una família
numerable en què cada \(A_n\) és enumerable. Aleshores la unió disjunta
\(\coprod_{n\in\N}A_n\) és enumerable, i per tant també ho és qualsevol
imatge seva ---en particular la unió indexada dins un conjunt ambient.
\end{teorema}

\begin{prova}
Per hipòtesi, per a cada \(n\) el conjunt \(A_n=p^{-1}(n)\) és enumerable:
o bé és inicial, o bé \emph{existeix} un epimorfisme
\(\N\twoheadrightarrow A_n\). Reunim totes les enumeracions de totes les
fibres en una sola part d'un objecte fix:
\[
E':=\bigl\{\,(n,e)\in\N\times E^{\N}\ :\ p\circ e=n\circ{!}_\N
\ \text{ i }\ \forall x\in E\,\bigl(p(x)=n\Rightarrow\exists k\,(e(k)=x)\bigr)\bigr\},
\]
tallada amb els quantificadors interns (capítol~\ref{cap:parts-interna});
la fibra de \(E'\to\N\) sobre \(n\) és el conjunt de les enumeracions de
\(A_n\), no buida exactament quan \(A_n\) no és inicial. Sigui
\(N'\hookrightarrow\N\) la part dels \(n\) amb \(A_n\) no inicial (la imatge
de \(p\)). Si \(N'\) és inicial, la unió és inicial i no hi ha res a
provar. Altrament, \(N'\) és enumerable sense elecció
(teorema~\ref{teo:subobjectes-N}): sigui \(r\colon\N\twoheadrightarrow N'\)
una enumeració. La família \(E'\) reindexada per \(r\) és una família
\(E''\twoheadrightarrow\N\) amb totes les fibres no buides, i l'axioma
d'elecció numerable en dona una secció: una \emph{tria simultània} d'una
enumeració \(e_{r(j)}\colon\N\twoheadrightarrow A_{r(j)}\) per a cada
\(j\). Per a cada \(n\in N'\) prenem \(e_n:=e_{r(j_n)}\), on \(j_n\) és
el \emph{mínim} índex \(j\) amb \(r(j)=n\), que existeix perquè \(r\) és
exhaustiva i és únic pel bon ordre de \(\N\): així l'assignació
\(n\mapsto e_n\) és explícita i no amaga cap elecció addicional. Amb aquesta tria, l'aplicació
\[
\N\times\N\longrightarrow\coprod_{n\in\N}A_n,
\qquad
(j,k)\longmapsto e_{r(j)}(k),
\]
és exhaustiva sobre punts (tot punt de la unió és a alguna fibra
\(A_{r(j)}\), i allà és algun \(e_{r(j)}(k)\)), per tant epimorfisme, i component-la amb un aparellament
\(\N\cong\N\times\N\) s'obté una enumeració de la unió disjunta. Tota
imatge d'aquesta unió és enumerable per la
proposició~\ref{prop:estabilitat-enumerable}(1).
\end{prova}

\begin{observacio}[On és l'elecció, i on no]\label{obs:eleccio-unio}
El pas irreductible és la tria simultània de les \(e_n\). Sabem que
cada \(A_n\) enumerable \emph{té} alguna enumeració, però la definició
d'enumerabilitat n'afirma només l'existència; res no en selecciona cap de
canònica quan els \(A_n\) són conjunts arbitraris. Encadenar aquestes
infinites eleccions en una sola aplicació \(n\mapsto e_n\) és precisament
el contingut de l'axioma~\ref{ax:eleccio-numerable}. Contrasta amb dos
casos ja vistos on l'elecció \emph{no} feia falta: la unió de \emph{dos}
enumerables (proposició~\ref{prop:estabilitat-enumerable}), on hi ha només
dues enumeracions donades; i el conjunt de llistes
(observació~\ref{obs:llistes-enumerables}), on les enumeracions dels
sumands eren totes derivades \emph{uniformement} d'una de sola. La
diferència no és la mida de la família ---en tots dos casos infinita en el
segon--- sinó si les enumeracions vénen donades o s'han de triar. De fet,
sense cap elecció és consistent que una unió numerable de conjunts de dos
elements no sigui enumerable: la conclusió del
teorema~\ref{teo:unio-numerable} pot fallar, i per tant l'axioma hi és
essencial.
\end{observacio}

\begin{observacio}[El producte numerable]\label{obs:producte-numerable}
El producte d'una família numerable, \(\prod_{n\in\N}A_n\), planteja una
qüestió d'elecció encara més aguda. El producte existeix com a objecte
---és un límit, i els límits numerables es poden construir a partir dels
exponencials i de \(\N\)---, però la seva \emph{no-vacuïtat} equival, per a
famílies de conjunts no buits, a una forma d'elecció: dir que
\(\prod_{n\in\N}A_n\) té un punt global és exactament triar un element de
cada \(A_n\). Per aquest producte l'elecció numerable és, doncs, no una
comoditat de la prova sinó la substància de l'enunciat. És el motiu pel
qual, en teories sense elecció, s'ha de distingir amb cura entre
l'existència d'un producte com a objecte i l'existència d'elements seus.
\end{observacio}

\begin{resumcap}
\apartat{Idees centrals}
\begin{enumerate}
  \item \emph{Infinit} és la negació d'\emph{finit}; \emph{infinit de Dedekind} és admetre un endomonomorfisme no exhaustiu. \(\N\) ho és pel successor; Dedekind implica infinit pel cas finit del capítol anterior, sense cap axioma.
  \item Infinit implica Dedekind amb l'elecció numerable, i només per un pas: triar alhora una llista injectiva de cada llargada. La successió injectiva es construeix després sense tria, prenent en llistes de llargada \(2^n\) el primer valor no usat.
  \item Enumerable és ser imatge de \(\N\); l'estabilitat per imatge, part, suma i producte és sense elecció, amb un aparellament \(\N\cong\N\times\N\) definit per recursió sobre les diagonals.
  \item Tota part de \(\N\) és enumerable, finita o isomorfa a \(\N\), perquè el mínim és una tria canònica; d'aquí, tot enumerable infinit és isomorf a \(\N\).
  \item \(\operatorname{List}(A)\) és una part d'\((A+\mathbf1)^{\N}\); el seu principi de recursió es demostra pujant a \(B^{\operatorname{List}(A)}\) i prenent la diagonal sobre la llargada. És enumerable, uniformement, si \(A\) ho és.
  \item La unió numerable d'enumerables necessita l'elecció numerable (i és consistent que falli sense ella); el producte numerable la incorpora al seu enunciat.
\end{enumerate}
\apartat{Errors típics} Confondre els dos usos de la paraula: \emph{enumerable} i \emph{numerable} qualifiquen conjunts (imatge de \(\N\); finit o enumerable infinit), mentre que \enquote{família, unió o producte numerable} i \enquote{elecció numerable} volen dir \emph{indexat per \(\N\)}; identificar infinit i infinit de Dedekind sense elecció; definir una successió \enquote{pel primer element no usat} com si fos una recursió de l'axioma (depèn de tota la història); creure que la unió de dos enumerables i la de numerablement molts estan al mateix nivell axiomàtic; pensar que \(\operatorname{List}(A)\) és \enquote{la suma \(\coprod_nA^{\seg n}\)}, que els axiomes no formen directament; i confondre l'existència del producte numerable com a objecte amb l'existència dels seus punts.
\apartat{Lectura recomanada} \cite[§§7.6--7.8]{lawvere-rosebrugh} per a Dedekind-infinitud, enumerabilitat i el paper de l'elecció; \cite[§§13, 15--16 i 23]{halmos} per als conjunts infinits, l'elecció i els conjunts numerables en la teoria clàssica; \cite[apèndix~B]{lawvere-rosebrugh} per a les formes de l'elecció; i \cite[pàg.~12]{bell2006} per a la definició de Dedekind i l'axioma d'infinitud com a existència d'un conjunt \(X\cong X+1\).
\end{resumcap}
